\documentclass[12pt]{article}
\usepackage[margin=1in]{geometry}
\usepackage{graphicx, url}
\title{Impacto da Inteligência Artificial no Mercado de Trabalho}
\author{Rael Fernando de Mello}
\date{\today}

\begin{document}
\maketitle

\section{Introdução}

    Nestes últimos anos, vêm se notando um grande crescimento das ferramentas de Inteligência Artificial. Estas tecnologias vieram com a função de melhorar o nosso cotidiano contribuindo para diversas situações, como por exemplo em pesquisas rápidas, explicações de conceitos,estudos e etc.

    Além disso, por possuirem um alto poder computacional, elas podem estar não apenas em áreas simples, mas também no mercado de trabalho, especialmente na área de TI. Com base nisso, neste presente artigo, iremos abordar principalmente as influências da tecnologia nesta área que com o advento da tecnologia vêm se ganhando bastante relevância e como ela pode influenciar positivamente ou negativamente os profissionais desta área.

    Com base nisso, será efatizado também sobre algumas características e sobre algumas adversidades sobre a mesma. Além do mais,também será criado um software na linguagem de programação Python com o auxílio de I.A para ser feita uma analise crítica sobre como ela influência no desenvolvimento de um programa.
\section{Fundamentação Teórica}
    \subsection{Aplicações}
        As ferramentas de Inteligência Artificial possuem inúmeras funções dentro delas que podem ser aplicados em diversas partes do cotidiano de um desenvolvedor. Algumas das principais aplicações da Inteligência Artificial na área de TI são:
        \begin{itemize}
            \item Relatórios: Um relatório de sobre os funcionários da empresa, setores, vendas, desenvolvimento de programas e entre outros assuntos, com apenas alguns cliques uma tarefa algumas horas ficou pronta em alguns segundos, poupando muito tempo, assim evitando o trabalho manual sobre.
            \item Desenvolvimento de um programa: O desenvolvimento de um programa como um site ou um aplicativo é uma tarefa que demanda bastante tempo, estudo e foco para se realizar.
            \item Segurança e detecção de problemas: As IAs são ensinadas para analisar padrões de comportamento de usuários e sistemas em tempo real. Com isso, alguma ação fora do padrão ou qualquer outra são dectadas automaticamente em milissegundos na maioria das vezes.  
        \end{itemize}

    \subsection{Benefícios}
        Como abordamos na introdução, a Inteligência Artificial possui, atualmente, um grande agregado ao mercado de TI, auxiliando no desenvolvimento de softwares através de sua alta capacidade computacional. Então vamos ver alguns benefícios no mercado de Ti.
        
        Otimização do Tempo por conta de seu alto poder computacional algumas tarefas são realizadas em milissegundos, como foi mencionado na emissão de relatórios de uma empresa de desenvolvimento web por exemplo. Pesquisa é outro lugar onde esta tecnologia aplimora muito a produção de informação.
            
        A I.A pode será a responsável por cortes custos para as empresas. Por exemplo, um funcionário realizando tarefas manualmente demanda tempo e  também gasto para a empresa. Entretando, com uma tecnologia dessas, as tarefas repetitivas são realizadas em menos tempo e também demandando menos gasto da empresa.
            
        Além do mais inteligências artificiais não sofrem de exaustão após inúmeras tarefas realizadas ou pelo excesso de informação. Dessa forma, a diferença de possíveis erros cometidos por uma inteligência artificial e um funcionário se tornam evidentes, pois mesmo que a I.A possa errar, será menos que um ser humano.

    \subsection{Malefícios}
        Como podemos ver, a Inteligência Artifical possui inúmeras aplicações e benefícios as nossas vidas, principalmente ao mercado de T.I auxiliando em desenvolvimento de programas e outrs. Todavia, esta tecnologia não possui apenas coisas boas em nosso cotidiano, podendo prejudicar em algumas partes e etc. 

        Um dos malefícios é a dependência, muito por conta de suas muitas funções, esta tecnologia pode atrair os olhos de muitas pessoas para aplicar muito em seu dia e até mesmo na área de T.I. Porém este uso elevado da ferramenta pode gerar depedência da ferramenta, prejudicando nossas habilidades cognitivas, criatividade e intuição principalmente no mercado de T.I onde por mais que a I.A realize o desejado para o dev, se o uso for excessivo não fará bem. Por isso, é necessário estabelecer um equilíbrio entre o uso de Inteligência Artificial e o pensamento intuitivo.

        A privacidade também vêm sendo afetado pelas tecnologias de I.A, pois normalmente recolhem e checam quantidades imensas de dados pessoas, assim podendo levantar discussões sobre este problema. Dessa forma, a coleta de dados utilizada também para treinamento destas ferramentas feri a privacidade de usuários, com seus dados sendo compartilhados com terceiros sem aviso previo.

        O aumento do desemprego também pode ser um fator de risco da I.A. As automações proporcionadas por estas ferramentas pode levar a perda de empregos em vários setores, com ênfase em funncionários pouco qualificados. Entretanto, veremos posteriormente que isso não é apenas um malefício, cm a I.A proporcionando a abertura de mais vagas.

    \subsection{Mercado de Trabalho e Competências}
        Notavelmente, com esta expectativa com a ascensão da I.A, o mercado de trabalho será alterado, principalmente acerca de vagas de empregos e opurtunidades na área de TI.Com isso, as empresas cada vez mais investem em I.A que em 2025 ultrapassou 30\% em relção ao ano anterior.Isso torna fica em evidÊncia, onde por conta de suas aplicações no trabalho, será necessário um aprendizado sobre como utiliza-las corretamente, exigindo um alto estudo.

        Além disso, se elevará a necessidade por profissionais do ramo, que possuam capacitação para utiliza-las com toda sua capacidade e obter ganhos e aprimorar divero setores do mercado de trabalho. Entretanto, toda essa qualificação exigirá diversas habilidades para que a I.A seja usada com toda capacidade computacional.

        A principal habilidade exigida é a engenharia de prompt, para o operador formular promps detalhados e claros pra um melhor entedimento da tecnologia, assim realizando o requerimento de maneira adequada. Além disso, será necessário compreender qual ferramenta de I.A será ideal para cada tarefa, verificando a especialidade da mesma.


    \section{Aplicação prática}

        Como mencionado anteriormente, desenvolveremos um programa bem simples, entretanto com auxílio de Inteligência Artificial para refletir sobre sua interferência no desenvolvimento de programas.
        Acerca do programa foi definido que ele será um sistema de cadastro de biblioteca, onde será cadastrado livros, e informações sobre os mesmos, será realizados empréstimos e outras opções sendo criado na linguagem Python.

        A 1ª versão do programa que foi criado sem auxilio de Inteligência Artificial contia dois arquivos: 
        \begin{itemize}
            \item Main.py : Onde o código era executado.
            \item Sistema.json : Onde os livros eram cadastrados.
        \end{itemize}
        Além disso, no Main.py havia as seguintes ações para o usuário:
        \begin{figure}[htbp]
            \centering
            \includegraphics[width=0.5\textwidth]{Ações Biblioteca.png}
            \caption{Ações}
            \label{fig:exemplo}
        \end{figure}
        Na primeria opção, os livros são cadastrados com nome, autor, data de Lançamento e número de páginas, além de avisar que estão disponíveis para empréstimos e salvos em Sistema.json.

        Na 2ª opção, é listado todos os livros já cadastrados, verificando se há livros para listar ou não. Na opção de Empréstimo, é informado qual livro deseja emprestar e então ele fica indisponível para empréstimo, verificando se existe livro que nem na anterior.
        Na 4ª opção é informado qual livro será excluído, passando pelas mesmas verificações anteriores, e enfim sendo deletado o livro.

        Após essa 1ª versão, foi enviada a Inteligência Artificial sendo o chatgpt a utilizada, um prompt contendo o contexto com o código escrito pedindo algumas sugetões de melhorias ao código, para melhora-lo em alguns sentidos. Segue o prompt abaixo:
        \begin{figure}[htbp]
            \centering
            \includegraphics[width=0.5\textwidth]{1ª Prompt.png}
            \caption{Prompt enviado}
            \label{fig:exemplo}
        \end{figure}
        Após isso, o Chat Gpt sugeriu algumas mudanças que serviam como um bom avanço para o programa desenvolvido, entre as principais melhorias estão:
        \begin{itemize}
            \item Funções para cada ação.
            \item Ação de Devolução.
            \item Validação Data de Lançamento.
            \item Aprimorar o sistema de Empréstimos.
        \end{itemize}
        
        \subsection{Código Final}

        O código final possui desta vez 3 arquivos, adicionando o Funções.py onde ficam todas as funções do programa incluindo as dos arquivos JSON e todas as ações do menu. O código continua com as mesmas ações incluindo a nova de devolução dos livros, além disso
        foi adicionado funções extras acerca da ação de empréstimo onde foi inplementado informações como usuário e data de devolução, ainda básicas mas bem formuladas. 

            Entre as principais melhorias, as que merecem destaque são:
        \begin{itemize}
            \item Ação de devolução: Agora é possível devolver o livro que anteriormente estava emprestado. Além disso, nesta ação também verificará se o livro está em empréstimo, exibindo se está tentando devolver um livro que não está em empréstimo. Entretanto, ainda não há verificação da data de devolução estar no limite, sendo uma limitação.

            \item Informações do empréstimo: Agora na função de empréstimo, é possível adicionar informações sobre quem pegou o livro e até que data o livro estará em empréstimo.

            \item Funções JSON: Foi inplementado também funções para selecionar e escrever no arquivo JSON onde possui o banco de dados dos livros, sendo uma reutilização de código. 

            \item Funções para cada ação: Para reaproveitar e organizar o código, foi criado o arquivo Funções.py o qual possui as funções do programa que realizam as ações anteriormente mencionadas, recebendo o parâmetro dados sendo o arquivo JSON convertido e aplicando a ação sobre ele.

            \item Validação da Data: Para a inserção de datas, foi utilizada a biblioteca datetime para uma validação ideal, sendo utilizada na data de lançamento do livro e data de devolução do livro.
        \end{itemize}

        
        

    \section{Analise dos resultados}

        Após inplementadas todas as outras melhorias, foi possível notar que programa obteve uma alta melhoria em seu funcionamento ficando mais organizado, onde desta vez possuia um arquivo nomeado Funções.py para guardar as funções do programa.

        Além do mais, o código agora possuia novas funções como de cadastrar informações do empréstimo do livro, contendo nome e data de devolução, mesmo que ainda bem simples mas que consolidavam uma base para programas mais complexos.

        Agora sobre a analise crítica, há algumas perguntas que devem ser espondidas para que possa notar a influência da I.A no desenvolvimento de um programa, entre elas estão.

        \subsection{O programa Funcionou corretamente ?}

        O programa antes e depois da melhoria proposta pela inteligência artificial funcionou corretamente pois foram corrigidos tabém alguns erros com auxilio naquele mesmo prompt.

        \subsection{Foi necessário corrigir o código ? Por quê ?}

        Foi necessário corrigir o código, pois faltava algumas partes importantes para um sistema que nem de uma biblioteca. Por exemplo, faltava a ação de devolução sendo uma ação bem importante. Não só mais também, foi colocado uma validação aprimorada para a data de lançamento sendo utilizada a biblioteca datetime nativa do Python.

        \subsection{Quais conhecimentos de programação foram necessários ?}

        É necessário para a construção desse programa a parte básica do python, englobando operadores aritméticos, estrutura condicional, laços de repetição, funções e etc. Além disso, foi necessário um conheciemnto mais aprofundado utilizando arquivos em formato JSON juntamente com a respectiva biblioteca e também utilizar módulos do Python.

        \subsection{Quais as vantagens e os risco do uso da IA ?}

        As principais vantagens da I.A no desenvolvimento web são a rápida solução de problemas, pois se emparcarmos a I.A pode solucionar tirando nossas dúvidas acerca do problema, e também de pedir sugestões de melhorias sobre o programa. Entretanto sobre as desvantagens estão principalmente como falamos anteriormente da depedência da ferramenta, prejudicando nossas criatovidade e pensamento cognitivo.

        \subsection{A I.A pode substituir os profissionais ou modificar suas atividades ?}

        Sobre ste assunto ainda é necessário acompanhar por um tempo sobre como o mercado de trabalho vai se acomodar com o advento da I.A, mas acredito que não irá substituir os profissionais de T.I no geral. ALgumas tarefas vão provavelmente ser susbtituídas, mas sendo apenas auxiliadas com a ferramenta. Além disso, também é necessário as pessoas que saibam usá-las adequadamente para assim um bom proveito da ferramenta.

        \subsection{Por que ainda é importante aprender programação}

        Ainda é importante aprender sobre programção pois é necessário o desenvolvedor entender sobre o que ele está escrevendo, como ele está realizando a tarefa, assim necessitando também de conhecimento deste ramo da computação. Não será necessário um largo conhecimento em programação, muito por conta da I.A que pode auxiliar em algumas etapas, mas um conhecimento básico é necessário.

    \section{Conclusão}

        Com base no que observamos ao longo do artigo, é evidente que a I.A fará mudanças no mercado de trabalho, mudando principalmente a maneira de como realizamos certas tarefas. Um bom exemplo é no desenvolvimento deste programa de biblioteca, onde a ferramenta de I.A auxiliou no desenvolvimento corrigindo alguns erros além da na inplementação de melhorias no sistema. 

        Com isso, podemos compreender que a Inteligência Artificial vai realmente nos auxiliar em diversas aplicações do cotidiano, melhorando na realização das mesmas. Entretanto, é necessário entender também sobre os problemas que a I.A pode produzir e busca a forma correta e segura de utilizar.

    \section{Referências}
    
     \begin{enumerate}
        \item opservices.
        \item blog.pareto.io.
        \item Forbes.
        \item Site Puc Pr.
        \item Alura.
        \item Site da Exame.com.
        \item Ferramenta de I.A Chat gpt para o desenvolvimento do Código.
     \end{enumerate}

        







\end{document}

